\section{被覆の構造：brane--brane と brane--image}\label{sec:cover}
被覆 U$(N_A)$ の自由エネルギー $f^{\cov}$ は A 型のデータだけで決まり、クロスキャップ部分 $g$ には入らない。しかし $E$ の源（$\tilde I$ で不変な状態）は被覆の brane--image のセクターにあり、$g$ の構造を読むには被覆の構造が要る。本節では被覆を brane と image に分け、brane--brane が \cite{KN17} の不変量そのものであること、brane--image が状態の空間の同定で理解できることを示す。

\subsection{brane と image の分割}\label{sec:split}
被覆を brane（$Y$）と image（$Y'$）に分けた生成関数を作ると、次が成り立つ（trefoil と $T(2,5)$ の 2 箱。全ブロックは\ref{app:cover}。3 箱は分けていない）。
\begin{itemize}
\item brane--brane のセクターは \cite{KN17} の U$(N_A)$ の不変量そのものである：一様な規格化の因子 $(t/q)^{\delta|R|/2}$ を除くと、$(t^{-1/2} - t^{1/2}) f^{\mathrm{bb}}_R$ は \cite{KN17} の t 側の $f^{\bar t}_R$（$a = \lambda^2 = t^{N_A}$、鏡像）に $R = [1], [2], [1,1]$ で全係数一致し、hatted（$p_1^2$ 成分を $z_q$、$p_2$ 成分を $q^{1/2} + q^{-1/2}$ で割る $M \otimes M$）は鏡像のため $R \leftrightarrow R^T$ で \cite{KN17} の $\hat f$ に一致する（trefoil と $T(2,5)$、tests/refined\_kauffman/test\_rk\_cover\_kn.py）。したがって $G$ 展開でき、全 $\beta$ で符号がそろう（\cite{KN17} の正値性）。
\item brane--image のセクター $x = W_{\mathrm{adj}} + \kappa - W_{\yone}^2$ は、$(t^{-1/2} - t^{1/2})x/z_q$ が整数 Laurent で $\sigma$ 対称である。$\lambda^\beta$ のブロックの多くは「そろう」か「$(\bt^2 - 1)\times$ そろう」かのどちらかである。混在するのは、trefoil では $\lambda^{-8}$ の 1 ブロック、$T(2,5)$ では中央の $\lambda^{-16}$〜$\lambda^{-10}$ の 4 ブロックである。
\end{itemize}
向き付けられた被覆の brane--image のセクターが、クロスキャップと同じ普遍因子 $(\bt^2 - 1)$ をもつ。
$x$ を $z_q$ で割る理由は命題 \ref{prop:rows} にある：$(1-t)x = s^2(t - \lambda^2)(\lambda^2 - 1)(q - 1)(\lambda^2 q - 1)\,N$ で、$(q-1)$ は brane と image の間の環状の弦の refined の測度 $(q-1)(t-1)$ の $q$ 側の因子、$(t-1)$ はモードのタワー $1/(1-t)$ と打ち消す（framed unknot の固定セクターの指数 $(\lambda^2-1)(q-1)\sum_k(\lambda^2q)^{\mathsf f-1-k}t^k$ も同じ因子をもつ）。brane と image の境界因子が非対称なのではない。
$\lambda$ の境界因子 $(t-\lambda^2)(\lambda^2-1)(\lambda^2q-1)$ は $\lambda^\beta$ のブロックごとに符号が交代するので、$x/z_q$ のブロックの混在の大半はこの因子による。境界因子をすべて外した $N$ を $G$ 展開すると（\ref{app:cover}）、ブロックは trefoil で全部そろい、$T(2,5)$・$T(2,7)$ では切り替わりのブロック $k = 2m$（$\lambda^{-(6+4m)}$）だけが混在する。そのブロックは $n_{2m} = \mathrm{Sym}^{2m}V - S\otimes S$（計算結果 \ref{cr:principle}）で、$\mathrm{Sym}^{2m}V$ と $S^2$ はそれぞれ $G$ 展開の係数がそろう（$m \le 3$）。したがって brane--image の混在は、対の状態の引き算 $-S\otimes S$ という状態の空間の構造そのものであり、composite の値 $R_{\mathrm{adj}}$ に固有の非正値性（§\ref{sec:composite}、\cite{CE15}）やクロスキャップの問題とは別である。
brane--image のセクターは、$\tilde I$ で自分に移る状態（$E$ の源）を含むセクターである。

\subsection{brane 1・image 1 のセクター：状態の空間の同定}\label{sec:principle}
固定セクターの最低次（brane 1・image 1、$n = 1$）では、$E$ は条件 (i)(ii)（§\ref{sec:hopf}）を使わずに状態の空間の同定だけから作れる。$T(2,2m+1)$ の $E$ の台 $N$ を $\lambda^{-2}$ の冪（ブロック $k = 0, \dots, K = 4m-1$）で分け、$V = 1 + t + qt$（trefoil のreduced superpolynomial）、$S$ を $T(2,2m+1)$ のreduced superpolynomial、$\psi^2$ を $(q, t)\to(q^2, t^2)$ とする。
% [P17]
\begin{compresult}[符号の規則の破れ]\label{cr:signrule}
下の系列（$k \le 2m$）に $(-1)^{\deg_q}$、上の系列に $(-1)^{\deg_t}$ を付ける符号の規則は、$T(2,3)$ と $T(2,5)$ では条件 (i)（$t = -1$）と (ii)（$t = q$）のもとで $E$ を一意に決めるが、$T(2,7)$ では (i) だけでも解が無い。破れるのは上の系列の最初のブロック $k = 2m+1$ である。
\end{compresult}
% [P18]
\begin{compresult}[状態の空間と $\tilde I$]\label{cr:principle}
$N$ のブロック $n_k$（$\lambda^{-(6+2k)}$ の係数を $r^{-(4+2k)}$ で正規化したもの）は
\[
n_k = \mathrm{Sym}^k V \ (k < 2m),\qquad n_{2m} = \mathrm{Sym}^{2m}V - S\otimes S,\qquad n_{K-j} = -t^{K-2j}\bigl[\mathrm{Sym}^j V + (q-t)\,R_j\bigr]\ (0 \le j \le 2m-2)
\]
であり、$\tilde I$ は $V$ の対合 $\mathrm{diag}(1, 1, -1)$（上の系列では双対）の誘導作用、$S\otimes S$ では因子の入れ替え（トレース $\psi^2 S$）として働く。この同定から作った $E$ は、$T(2,3)$・$T(2,5)$ で (i)+(ii) から決まる一意な $E$ を全単項式で再現し、$T(2,7)$〜$T(2,15)$ では予言として (i)・(ii) をともに満たす。
境界の因子を外した $N = (1-t)x/(s^2(t-\lambda^2)(\lambda^2-1)(q-1)(\lambda^2q-1))$ と $\tilde N = E/(s^2(t-\lambda^2)(\lambda^2+1)(q+1)(\lambda^2q+1))$ を単項式ごとに比べると、$T(2,3)$〜$T(2,15)$ の全単項式で $|\tilde e_m| \le |c_m|$、$\tilde e_m \equiv c_m \pmod 2$ が成り立ち、$k < 2m$ のブロックでは $|\tilde e_m| = |c_m|$（全状態が $\tilde I$ で不変）、$k \ge 2m$ で初めて $\tilde I$ で入れ替わる対が現れる。$E$ は状態の数の符号付き部分和として読める。
\end{compresult}
% [P19]
\begin{conjecture}[補正の形]\label{conj:correction}
$R_{m+j'} = t^{m-1-j'}F_{j'}$ で、$F_{j'}$ は $m$ によらず
\[
\sum_{j'\ge 0}F_{j'}z^{j'} = \frac{1 + (2t + t^2 + qt + 2qt^2)z + qt^3z^2}{(1-z)(1-qtz)(1-t^2z)(1-q^2t^2z)}
\]
（状態の数は $\binom{2j'+3}{3}$）。$j' \le 5$（$T(2,7)$〜$T(2,15)$）で全係数が一致する。
\end{conjecture}
% [P20]
\begin{conjecture}[補正への $\tilde I$ の作用]\label{conj:correctionI}
$E$ の補正は $(-1)^{m+j'+1}t^{m-1-j'}\psi^2(X_{j'})$、$X_{j'} = \mathrm{Sym}^{j'}V + t\,\mathrm{Sym}^{j'-1}V$（母関数 $(1+tz)/((1-z)(1-tz)(1-qtz))$、$(j'+1)^2$ 状態）であり、$F_{j'}$ の残りの状態は $\tilde I$ で入れ替わる対である：$F_{j'} = \psi^2 X_{j'} + 2P_{j'}$、$\sum P_{j'}z^{j'} = tz(V' + qt^2z)/((1-z)(1-qtz)(1-t^2z)(1-q^2t^2z))$、$V' = 1 + q + qt$。$j' = 2$ は $T(2,9)$ で (i) の行と台から 32 通りの変種のうち一意に選ばれ、$j' = 3, 4, 5$ は $T(2,11)$, $T(2,13)$, $T(2,15)$ で (i)+(ii) と整合する。
\end{conjecture}
% [P23]
\begin{conjecture}[閉じた式]\label{conj:closedform}
$T(2,2m+1)$ の $n = 1$ の $N$ と $E$ は、結び目によらない普遍的なデータ $V$、$F_{j'}$、$X_{j'}$ と、結び目のデータ $S_m$（reduced superpolynomial）・$m$ だけで書ける（$K = 4m-1$、$n_k$ は $\lambda^{-(6+2k)}q^{-(2+k)}$ の係数）：
\[
 n_k = \mathrm{Sym}^kV\ (k<2m),\quad n_{2m} = \mathrm{Sym}^{2m}V - S_m^2,\quad n_{K-j} = -t^{K-2j}\bigl[\mathrm{Sym}^jV + (q-t)\,t^{m-1-j'}F_{j'}\bigr]\ (j = m+j' \ge m),
\]
$E$ は対応する $\tilde I$ のトレースである。$T(2,3)$〜$T(2,15)$ で $N$ は全係数で成立し、$E$ は (i)・(ii) を満たす（計算）。
\end{conjecture}
% [P24]
\begin{compresult}\label{cr:plog}
母関数の plethystic log は、$X_{j'}$ では有限（$(1+2t+qt)z - t^2z^2$：$V\oplus t$ の Fock 空間に関係 $t^2z^2$）だが、$F_{j'}$ と上の系列 $\mathrm{Sym}^jV + (q-t)t^{m-1-j'}F_{j'}$ では無限で、上の系列の生成元は $Vz$ と $(q-t)t^{m-1-j'}c_{j'}z^{m+j'}$（$c_0 = 1$、$c_1 = V^2 - t(1+qt)$、$c_2$ から負の係数）である。上の系列は正の BPS 生成元の Fock 空間としては読めない。
\end{compresult}
% [P25]
\begin{compresult}\label{cr:t34}
$n = 1$ の原理のうち、下の系列 $\mathrm{Sym}^kV$、切り替わり（普遍的な系列 $\mathrm{Sym}^jV$ から結び目の部分へ移る $j = 2\delta$ の項。以下「切り替わり」）で完全平方を引く構造、$K = rs - 3$、条件 (i)(ii) を満たす $E$ の存在は、2 ストランドでない $T(3,4)$（$\gamma = \tau_+\tau_-^3$、$K = 9$）でも成り立つ。しかし切り替わりの平方の因子は $\tilde S = (1+qt)(1+t+t^2+qt^2+q^2t^2)$（10 状態）で $T(3,4)$ のreduced superpolynomial（11 状態）ではなく、上の系列は $-t^{K-2j}[\mathrm{Sym}^jV + (q-t)R_j]$ の形にならない（$j = 0$ で $-t^8(q+t)$）。下の系列を固定した整数計画は (i)+(ii) で解をもつが一意でない（111 単項式中 29 が不定）。切り替わりの $S\otimes S$ と補正 $F_{j'}$・$X_{j'}$ の構造は 2 ストランドのトーラス結び目に固有である。
\end{compresult}
% [P26]
\begin{proposition}\label{prop:rows}
$x = \lambda^{-2}$ とし、$[1]$ の鏡像の $a$ の行を $R_1 = q^{-\delta/2}t^{-\delta/2}x^\delta R(x)$、$R(y) = \sum_j \rho_j y^j$、随伴 (composite $(1;1)$) の行を $R_{\mathrm{adj}} = q^{-2\delta}x^{2\delta}A(x)$、$A(y) = \sum_j \alpha_j y^j$ と書く。$n = 1$ の固定セクター $N = \sum_k n_k x^k$ は
\[
\sum_k n_k x^k = \frac{1 - x^{2\delta}R(qx)^2}{(1-x)(1-tx)(1-qtx)} + \frac{x^{2\delta-1}G(qx)}{1-tx},\qquad A = R^2 + (q-1)(t-1)G
\]
を満たし、$N$ が Laurent であることと $G$ が Laurent であることは同値である。
\end{proposition}
\begin{compresult}\label{cr:rows}
$T(2,3)$〜$T(2,15)$、$T(3,4)$、$T(3,5)$、$T(3,7)$、$T(3,8)$、$T(3,10)$、$T(4,5)$、$T(4,7)$ の 14 のトーラス結び目で $g_0 = 0$、$g_1 = (2\rho_0\rho_1 - \rho_1^2)/q$ が成り立ち、したがって切り替わりは $n_{2\delta} = \mathrm{Sym}^{2\delta}V - (\rho_0 - \rho_1)^2$ である。平方の因子は常に 表現 [1] の最上 2 行 $\rho_0 - \rho_1$ であり、2 ストランドではこれがreduced superpolynomial に一致し（行が 2 つ）、$T(3,4)$ の $\tilde S$ は 3 行目 $t^3$ を切り捨てたものである。$N$ は随伴と表現 [1] の被覆データだけで決まるので、$T(3,5)$（$\gamma = \tau_+\tau_-\tau_+\tau_-$、$\delta = 4$、$K = 12$）、$T(3,7)$、$T(3,8)$、$T(3,10)$、$T(4,5)$、$T(4,7)$ では $\tilde S$ と $K = rs - 3$ を予言として的中させた。最深ブロックは $n_K = -t^{K-s+2}h_{s-2}(q,t)$（$s = 2, 3, 4$）である。
\end{compresult}
% [P27, P28]
\begin{compresult}\label{cr:threestrand}
3 ストランドの上の系列 $U_j = -n_{K-j}/t^{K-2j}$ は $j \le \delta - 2$ で $(1 + q/t)\mathrm{Sym}^jV + (q-t)C_j$ の形をもち（$C_j$ は結び目に依る）、$j = \delta - 1$ で $t = q$ の完全平方 $-([\lfloor r/2\rfloor]_{t^2}[r']_t)^2/t^{\delta-1}$（$r'$ は $r$ 以下の最大の奇数）が入る（$T(3,4)$、$T(3,5)$、$T(3,7)$、$T(3,8)$、$T(3,10)$）。4 ストランド（$T(4,5)$、$T(4,7)$）では前因子が $h_2(q,t)/t^2$、平方の位置が $j = r - s + 1$ で、平方の因子は $\rho_0(T(s-1, r-s+1))\cdot S(T(2,r'))|_{t=q}$ と読める。$E$ については、切り替わりの規則 $E_{2\delta} = \mathrm{Sym}^{2\delta}V|_{\varepsilon} - \psi^2\tilde S$ が $T(3,4)$・$T(3,5)$ でも条件 (i)(ii) と両立し、最深の 2 ブロックが一意に決まる（$E_K = -t^8(q-t)$、$+t^{11}(q-t)$）。2 ストランドでは、$E$ を同じ $(R, G)$ の形に分けると $G$ の行の $\tilde I$ は深い行（$j \ge \delta + 2$）で $q \mapsto -q$ の規則、それより上で対の補正 $D_{\delta-i} = t^{a(m,i)}\tilde D_i$（$\tilde D_i$ は $i$ と $m$ の偶奇だけで決まる）である。
\end{compresult}
% [P29]
\begin{compresult}[composite の大きな $N$ の因子化]\label{cr:composite}
命題 \ref{prop:rows} の $A = R^2 + (q-1)(t-1)G$ は随伴に限らない。trefoil の composite $(2;1)$、$(1,1;1)$、$(2;2)$、$(2;1,1)$、$(1,1;1,1)$ の大きな $N$ の値は
$R_{[\alpha;\beta]} = R_\alpha R_\beta + (q-1)(t-1)G_{\alpha\beta}$（最上行は積、差は全行で $(q-1)(t-1)$ で割れる）を満たし、$G_{\alpha\beta}$ の行は $\alpha, \beta$ に応じた次元型の因子 $(1+q)$、$(1+t)$ を含む。$(q-1)(t-1)$ は brane と image の間の環状の弦の refined の測度で、$N$ の整数性は連結相関の LMOV 型の整数性と同値である（読み）。$g_1 = (2\rho_0\rho_1 - \rho_1^2)/q$ の一般化 $\rho_0^\alpha\rho_1^\beta + \rho_1^\alpha\rho_0^\beta - \rho_1^\alpha\rho_1^\beta$ は他の composite では成り立たない。
\end{compresult}
% [P21]
\begin{proposition}\label{prop:pairblind}
$(q-t)$ の補正の対 $(q\mu, -t\mu)$ の $\tilde I$ の符号は条件 (ii) に現れない（$t = q$ で消える）。したがって補正の符号は (i) だけからは決まらず（$n = 1$ の $j' = 2$ で 15 中 2 個、$n = 2$ の $\lambda = (2)$ で 17 中 3 個）、予想 \ref{conj:correctionI} のような状態の空間の構造が要る。
\end{proposition}
したがって $n = 1$ では、$E$ と $C$ の分離（命題 \ref{prop:EC}）は状態の空間の同定で解消する。$n = 2$（brane 2・image 2）では、$N_\lambda$ が Wilson loop の無い背景（6 セルの Fock 空間、閉じた形）と結び目の部分に分かれ、$\lambda^{-12}$ までは「背景のセルの符号 $+$ 平方と行の積に入れ替え」で両 $R$ の $E$ が符号まで決まる（計算結果 \ref{cr:n2E}）。$\lambda^{-14}$ 以降は、$\sigma$ 共変な部分を brane の対象の平方の和として書く原理が無く、$E$ は未決定である（計算結果 \ref{cr:n2negative}、\S\ref{sec:discussion}）。

